\documentclass[12pt, a4paper]{article}
\usepackage[utf8]{inputenc}
\usepackage[T1]{fontenc}
\usepackage[slovak]{babel}
\usepackage{geometry}

\title{Metódy inžinierskej práce}
\author{Illia Kukhmanskyi, Kyryll Krasnikov, Oleksii Kovalevskyi, Nikita Kushchenko}
\date{\today}

\begin{document}

\maketitle

\section*{Spresnenie zamerania projektu}

\noindent\textbf{Názov projektu:} \\
Inteligentný asistent pre ochranu súkromia a správu digitálnej stopy

\vspace{0.5cm}
\noindent\textbf{Akronym projektu:} \\
PrivAs (Privacy Assistant)

\vspace{0.5cm}
\noindent\textbf{Zameranie projektu:} \\
Vývoj webovej platformy v oblasti Cybersecurity, Data Privacy a GDPR Compliance pre automatizovaný monitoring, správu a odvolávanie súhlasov so spracovaním osobných údajov vo webových službách.

\vspace{0.5cm}
\noindent\textbf{Anotácia projektu:} \\
Predkladaný projekt sa zameriava na návrh a implementáciu inteligentného asistenta pre ochranu súkromia a správu digitálnej stopy používateľov. V súčasnosti používatelia pri interakcii s digitálnymi službami často bez detailného prečítania súhlasia s rozsiahlymi podmienkami používania (Terms of Service), čím strácajú kontrolu nad svojimi osobnými údajmi. Zrušenie účtu navyše spravidla neznamená automatické vymazanie týchto dát z databáz poskytovateľa, pričom proces uplatnenia práva na zabudnutie (podľa legislatívy GDPR) je pre bežného používateľa časovo a administratívne náročný.

Cieľom projektu je vytvoriť softvérové riešenie (vo forme interaktívneho bota), ktoré pomocou technológií spracovania prirodzeného jazyka automaticky analyzuje zložité právne texty a extrahuje z nich kľúčové informácie o spracovaní dát či skrytých rizikách (napríklad predaj údajov tretím stranám). Systém bude zároveň slúžiť ako centralizovaný register (Consent Hub) udelených súhlasov a umožní ich automatizované spravovanie. V prípade požiadavky na výmaz údajov systém vygeneruje a odošle právne záväznú žiadosť podpore poskytovateľa služby v zmysle príslušných článkov GDPR (najmä článku 17).

\end{document}